\section*{Chapter \ref{ch:b1:economic-data-python} --- Economic Data with Python}
% Numbers recomputed in code/bachelor-1/28-economic-data-python/exercises/ (solutions.py, exo_05.py,
% exo_06.py, exo_09.py, exo_10.py, exo_11.py; tested in tests/test_exercises.py).

\begin{solution}{exo:b1:economic-data-python:shapes-ex}
First look: 18 rows and 6 columns, a country in a year, key \texttt{(iso3, year)}.
Long table: 25 rows and 3 columns, same unit and key; wide table: 5 rows
(countries) and 5 columns (years), key \texttt{iso3}. Monthly series: one value per
month, key \texttt{(year, month)}.
\end{solution}

\begin{solution}{exo:b1:economic-data-python:yoy}
$9.06\%$ unadjusted, $8.98\%$ adjusted. Over twelve months the season cancels
almost exactly; the small difference comes from the adjusted series'
factors changing slowly from year to year.
\end{solution}

\begin{solution}{exo:b1:economic-data-python:annualise}
$1.004^{12} - 1 = 4.91\%$.
\end{solution}

\begin{solution}{exo:b1:economic-data-python:aggregate}
Unemployment rate: mean. Exports (a flow): sum. Consumer price index: mean.
Population: mean for a yearly average, or the last value for the end of
the year, said which. Debt (a stock): last value.
\end{solution}

\begin{solution}{exo:b1:economic-data-python:maddison}
161 keys match; 24 are only in the Penn table and 8 only in the Maddison
table. Only in Penn: small economies such as Aruba, Anguilla and the Cayman
Islands. Only in Maddison: units that no longer exist (the Soviet Union,
Yugoslavia, Czechoslovakia, kept as historical aggregates) and economies
Penn does not cover (Afghanistan, Cuba).
\end{solution}

\begin{solution}{exo:b1:economic-data-python:rebase}
$69.7$ in 2015, $77.8$ in 2020 and $100$ in 2025; mean log growth $3.61$ log points a
year.
\end{solution}

\begin{solution}{exo:b1:economic-data-python:toy}
The centred moving average of $100 + t + s_t$ is exactly $100 + t$, since the
season sums to zero over four quarters. The differences are then $s_t$, and
their averages by quarter are $+2$, $-1$, $-2$ and $+1$.
\end{solution}

\begin{solution}{exo:b1:economic-data-python:falling}
Positive: GDP per person grows faster than GDP, by about the rate at which
population falls.
\end{solution}

\begin{solution}{exo:b1:economic-data-python:hook}
$2.90\%$ nominal, $1.28\%$ in volumes, $0.98\%$ per person. The nominal rate changes
most (from $3.24\%$ to $2.90\%$): the inflation of 2021--23 is left out.
\end{solution}

\begin{solution}{exo:b1:economic-data-python:quarterly}
As \texttt{adjust}, with a four-quarter moving average centred with weights $1/8$,
$1/4$, $1/4$, $1/4$, $1/8$ and factors averaged by quarter. On the known series it
recovers $0.98$, $1.0099$, $1.03$ and $0.9801$: the season to within $0.0001$.
\end{solution}

\begin{solution}{exo:b1:economic-data-python:silent}
All 199 rows of the World Bank table, none with Penn data: no code matches,
and a left join keeps the rows with empty columns. Counting the matches
(\texttt{indicator=True}) or checking that the joined columns are not all empty
would catch it; after converting the codes to upper case, 178 rows match.
\end{solution}

\begin{solution}{exo:b1:economic-data-python:base}
In the stated base year, real and nominal GDP must be equal. If they are
not, the deflator's base year is not the one stated.
\end{solution}

\begin{solution}{pb:b1:economic-data-python:portal}
\textbf{1.} 26 rows, 2000 to 2025.

\textbf{2.} The largest absolute difference is 0: the identity holds.

\textbf{3.} The largest difference is 0.

\textbf{4.} 26 and 130: five goods in each of 26 years.

\textbf{5.} 100: the index is based in 2015, as stated.

\textbf{6.} A table that fails its own identities makes every later result
wrong; it is cheaper to find that out first.

\textbf{7.} 2021.

\textbf{8.} 2012: the 2021 row was typed as 2012.

\textbf{9.} Raised an error, because 2012 appears twice in the trade file.

\textbf{10.} With 2021 restored, all 26 years match one to one.

\textbf{11.} $40{,}231/129{,}777 = 31.0\%$.

\textbf{12.} $40{,}231/1.242 = 32{,}392$ million kests.

\textbf{13.} $1 + 0.004\sin(\pi/2) = 1.004$.

\textbf{14.} $1.0041$.

\textbf{15.} $0.996$ true and recovered.

\textbf{16.} $0.0003$.

\textbf{17.} $3.19\%$ over the twelve months to December 2024, against $3.79\%$ for the
yearly average.

\textbf{18.} The yearly average for a year's inflation, as the statistical
office publishes it; the December change for where inflation stood at the
end of the year. Either, if the report says which.

\textbf{19.} From nominal GDP: $(393{,}264/248{,}438)^{1/10} - 1 = 4.70\%$, which mixes price and
volume.

\textbf{20.} $1.64\%$ a year.

\textbf{21.} $1.63$ log points: $\ln(1 + g) < g$ for $g > 0$.

\textbf{22.} $1.04\%$ a year.

\textbf{23.} $0.99\%$: the consumer price index rose faster than the GDP
deflator ($43.5\%$ against $34.5\%$), since it prices a different basket.

\textbf{24.} The corrected \omterm{prop:b1:models-and-data:real-growth}{real growth} of Brenland, 2015--2025, is $1.6\%$ a year, not
the portal's $4.7\%$, which was nominal growth.
\end{solution}
